\section{Background}\label{sec:background}

\subsection{DKIM}
DomainKeys Identified Mail~\cite{rfc6376} lets the organisation that owns a
signing domain take responsibility for a message. The signer adds a
\texttt{DKIM-Signature} header field naming the signing domain
(\texttt{d=}), a selector (\texttt{s=}) under which the public key is
published in DNS at \texttt{s.\_domainkey.d}, the list of signed header
fields (\texttt{h=}), a hash of the canonicalized body (\texttt{bh=}), the
canonicalization (\texttt{c=}) and an optional body length limit
(\texttt{l=}). The signature covers the selected header fields, each taken
from the bottom of the header block upward, followed by the
\texttt{DKIM-Signature} field itself with its \texttt{b=} value empty. Under
\emph{simple} header canonicalization the fields are signed as they are;
under \emph{relaxed} canonicalization names are lowercased, folded lines are
unfolded, and runs of whitespace become one space. Naming a field in
\texttt{h=} more often than it occurs (``oversigning'') makes later
additions of that field break the signature. RSA keys have at least 1024
bits~\cite{rfc8301}. A DKIM signature attests to the signing domain, not to
the author: DMARC~\cite{rfc7489} calls a signature \emph{aligned} when
\texttt{d=} matches the domain of the \texttt{From} address, whose syntax,
with display names, quoting, comments and folding, is
RFC~5322's~\cite{rfc5322}.

\subsection{Smart-account recovery on Perch}
Soroban is the smart-contract platform of the Stellar
network~\cite{soroban-docs}. A contract account authorizes through its own
\texttt{\_\_check\_auth}, which the host calls when a contract it is invoked
by requires its authorization; the call receives the signature payload, the
account's signature value from the transaction's authorization entry, and
the contexts being authorized. Perch~\cite{perch-repo} builds smart accounts
on OpenZeppelin's Stellar contracts~\cite{oz-stellar-contracts} and specifies
their recovery~\cite{perch-recovery-spec}.

An account's policy document may carry a \emph{recovery member}: a profile,
a mode, a controller, the replaceable signers, and timing in ledgers. Under
the \texttt{Loss} profile the owner key may be lost but not stolen, so the
owner can veto any recovery; under \texttt{Protected} it may be stolen, so
changes to recovery also need the recovery condition. A
\texttt{guardian-only} mode lists guardians and a quorum, \texttt{zk-only}
enrolls one zero-knowledge credential, and \texttt{combined} needs both.
Perch binds the member, guardian addresses included, into a
\texttt{config\_hash}.

Anyone may open a recovery attempt, which grants nothing by itself. Each
attempt binds one \emph{recovery statement}, a fixed-width encoding of the
action (lost-key, compromise, cancel, reconfigure or upgrade), the network,
the account, the controller, the configuration epoch and hash, timing, and
an action-specific subject; for a lost-key recovery the subject names the
attempt id and commits to the replacement credentials and the target
document. Guardians authorize the statement's digest
$\digest = \Sha(\text{statement})$ through the controller's
\texttt{submit\_guardian}, which calls
\texttt{guardian.require\_auth\_for\_args((d,))} and counts each guardian at
most once per attempt and evidence domain (initiation or cancellation);
reconfigurations and upgrades collect approvals of a change digest the same
way through \texttt{approve\_change}. When the quorum is reached the attempt is authorized;
it completes only after the delay, during which a \texttt{Loss} owner may
cancel it. Every configuration change and completion bumps the epoch, which
invalidates older attempts and evidence.

\subsection{Noir, UltraHonk and Soroban verification}
Barretenberg~\cite{barretenberg-docs} proves Noir~\cite{noir-docs} circuits
with UltraHonk, a Plonk-family~\cite{plonk} proof system whose verifier runs
a sumcheck, as in HyperPlonk~\cite{hyperplonk}, under a universal reference
string, so a new circuit needs a new verification key but no setup of its
own. Perch verifies UltraHonk proofs in Barretenberg's zero-knowledge flavor
with a Keccak transcript in a Soroban contract; proofs are 16{,}224 bytes
whatever the circuit's size\claim{circuit-size}, and the email circuit uses
the same verifier code with its own key. Poseidon2~\cite{poseidon2} is the
hash shared by Perch's circuit and Soroban's host.
